% =====================================================
% Plantilla adaptada de FloripaSat-2 / GOLDS-UFSC Cap. 9 (AIT Plan)
% Para usar en: 00_documento_principal/chapters/11_verificacion.tex (§11.2, §11.3)
% =====================================================

\section{Instrucciones de Ensamblaje y FlatSat}

\subsection{FlatSat (Pruebas sobre mesa)}
La configuración \textit{FlatSat} consiste en interconectar todas las placas del satélite
fuera del chasis, mediante cables puente, para realizar pruebas
\textit{Hardware-In-the-Loop} (HIL) e identificar problemas de interfaz antes del
ensamblaje mecánico final.

\textbf{Secuencia FlatSat:}
\begin{enumerate}
    \item Conexión de placas en mesa según topología PC-104.
    \item Verificación de continuidad y rieles de alimentación (3.3 V, 5 V, vBatt).
    \item Encendido controlado con fuente de laboratorio (current-limited).
    \item Verificación de \textit{boot} y handshake OBC$\leftrightarrow$EPS$\leftrightarrow$UHF.
    \item Ejecución de comandos típicos (housekeeping, beacon, payload trigger).
    \item Validación de modos de operación de la OBC (Inicio, Rutina, Seguridad, Decomisión).
\end{enumerate}

\subsection{Procedimiento de Integración Mecánica}
\begin{enumerate}
    \item Limpieza de superficies y verificación dimensional de la estructura.
    \item Montaje de placas en el bus PC-104 con torques estandarizados (TBD~N·m por tornillo M3).
    \item Conexión del arnés interno según diagrama de cableado.
    \item Montaje de paneles solares y antenas plegadas.
    \item Inserción de kill-switches y verificación de continuidad RBF (Remove-Before-Flight).
    \item Cierre estructural final y verificación de envolvente.
\end{enumerate}

\section{Plan de pruebas físicas y validación ambiental}

\subsection{Pruebas Físico-Mecánicas}
\begin{itemize}
    \item \textbf{Dimensiones (Dimensions):} medición CMM/scanner contra envolvente CubeSat 2U.
    Tolerancia: $\pm 0{,}1$~mm en rieles.

    \item \textbf{Verificación de interfaz física (Fit Check):} inserción y extracción dentro
    del P-POD inyector. Verificar holguras y orientación de kill-switches.

    \item \textbf{Masa y Centro de Gravedad (CG):} balanza de precisión y método de
    suspensión. CG debe ubicarse a $< 2{,}0$~cm del centro geométrico en los 3 ejes.

    \item \textbf{Vibración Sinusoidal y Aleatoria (Random):} mesa vibratoria, perfil
    representativo del lanzador (TBD: Falcon 9, PSLV, Soyuz, etc.) en los 3 ejes (X, Y, Z).
    \begin{itemize}
        \item Sweep sinusoidal: 5--100~Hz, $\pm 2$~mm o 2.5~g.
        \item Random: perfil \textit{Component Minimum Workmanship} del estándar GSFC-STD-7000.
    \end{itemize}

    \item \textbf{Shock:} medio sinusoidal o pirotécnico simulado, 1000~g pico, 0.5~ms.
\end{itemize}

\subsection{Pruebas de Ambiente Espacial}
\begin{itemize}
    \item \textbf{Ciclos Térmicos (Thermal Cycling):} cámara $-15$ a $+60\,^\circ$C, 10
    ciclos. El OBC debe seguir transmitiendo vía umbilical. Detección de termofatiga
    en soldaduras críticas.

    \item \textbf{Vacío Térmico (TVAC + Bake-Out):} $+60\,^\circ$C en alto vacío durante
    semanas para promover desgasificación (outgassing) y certificar limpieza para
    integración en cohete.

    \item \textbf{Compatibilidad Electromagnética (EMC):} prueba de inmunidad e intermodulación
    entre subsistemas. Especial énfasis en el receptor UHF/VHF para descartar autointerferencia.
\end{itemize}

\subsection{Secuencia recomendada}
\begin{enumerate}
    \item Inspección visual y verificación de masa/CG.
    \item Fit check en P-POD.
    \item Vibración (X, Y, Z secuencial).
    \item Shock.
    \item Re-inspección visual y eléctrica.
    \item Ciclos térmicos.
    \item TVAC + Bake-Out.
    \item EMC.
    \item Test funcional final (FlatSat-equivalente sobre satélite integrado).
\end{enumerate}
